%%%%%%%%%%  scan idx 1  |  folio cover-verso  %%%%%%%%%%
% The contemporary-template restitution is documented in the new edition preface and source records.

% The masked notice is restored from the contemporary primary template
% in Springer LNM 228 (1971), front-matter PDF page 1; see the separate
% ASTRA-FRONT-NOTICE-1971-RESTORATION evidence packet. No obscured stroke
% in the supplied SGA 6 image is claimed to have become directly readable.
This series aims to report new developments in mathematical research and teaching -- quickly, informally
and at a high level. The type of material considered for publication includes:
\begin{enumerate}
\item Preliminary drafts of original papers and monographs
\item Lectures on a new field, or presenting a new angle on a classical field
\item Seminar work-outs
\item Reports of meetings
\end{enumerate}
Texts which are out of print but still in demand may also be considered if they fall within these categories.

The timeliness of a manuscript is more important than its form, which may be unfinished or tentative.
Thus, in some instances, proofs may be merely outlined and results presented which have been or will
later be published elsewhere.

Manuscripts should comprise not less than 100 pages.

Publication of \textit{Lecture Notes} is intended as a service to the international mathematical community, in
that a commercial publisher, Springer-Verlag, can offer a wider distribution to documents which would
otherwise have a restricted readership. Once published and copyrighted, they can be documented in the
scientific literature. --

\subsection*{Manuscripts}
Manuscripts are reproduced by a photographic process; they must therefore be typed with extreme care.

Symbols not on the typewriter should be inserted by hand in indelible black ink.
Corrections to the typescript should be made by sticking the amended text over
the old one, or by obliterating errors with white correcting fluid. Authors
receive 75 free copies.

The typescript is reduced slightly in size during reproduction; best results
will not be obtained unless the text on any one page is kept within the overall
limit of $18\times26.5$ cm ($7\times10\frac12$ inches). The publishers will be
pleased to supply on request special stationery with the typing area outlined.

Manuscripts in English, German or French should be sent to Prof. Dr. A. Dold,
Mathematisches Institut der Universität Heidelberg, 69 Heidelberg/Germany,
Tiergartenstraße or Prof. Dr. B. Eckmann, Eidgenössische Technische Hochschule,
CH-8006 Zürich/Switzerland.

\subsection*{Lecture Notes in Physics}
\textbf{Bisher erschienen/Already published}

Vol. 1: J. C. Erdmann, Wärmeleitung in Kristallen, theoretische Grundlagen und
fortgeschrittene experimentelle Methoden. II, 283 Seiten. 1969. DM 20,--

Vol. 2: K. Hepp, Théorie de la renormalisation. III, 215 pages. 1969. DM 18,--

Vol. 3: A. Martin, Scattering Theory: Unitarity, Analytic and Crossing.
IV, 125 pages. 1969. DM 14,--

Vol. 4: G. Ludwig, Deutung des Begriffs physikalische Theorie und axiomatische
Grundlegung der Hilbertraumstruktur der Quantenmechanik durch Hauptsätze des
Messens. XI, 469 Seiten. 1970. DM 28,--

Vol. 5: M. Schaaf, The Reduction of the Product of Two Irreducible Unitary
Representations of the Proper Orthochronous Quantummechanical Poincaré Group.
IV, 120 pages. 1970. DM 14,--

Vol. 6: Group Representations in Mathematics and Physics. Edited by V. Bargmann.
V, 340 pages. 1970. DM 24,--

Vol. 7: R. Balescu, J. L. Lebowitz, I. Prigogine, P. Résibois, Z. W. Salsburg,
Lectures in Statistical Physics. V, 181 pages. 1971. DM 18,--

Vol. 8: Proceedings of the Second International Conference on Numerical Methods
in Fluid Dynamics. Edited by M. Holt. IX, 462 pages. 1971. DM 28,--

